%%% FIELD ass-uniform-poly-time-model-condition
Fix the following uniform two-carrier computational convention.  First, a
causal instance \(G=(V,D,B)\), together with its supplied causal-vertex
order \(\triangleleft_V\), is relabeled by the order isomorphism
\(V\simeq\mathrm{Fin}\,|V|\) and encoded by its directed and bidirected
adjacency bit-matrices, the index of \(s\), and the binary numerators and
denominators of the positive rational costs; write \(L\) for the total
binary encoding length.  Second, and independently, a source graph
\(H=(U,E)\), together with its supplied source order \(\prec\), is relabeled
by the order isomorphism \(U\simeq\mathrm{Fin}\,|U|\) and encoded by its
adjacency bit-matrix and any accompanying rational data in binary, again
with \(L\) denoting the applicable total input length.  A polynomial-time
claim means polynomial time in the relevant carrier size and \(L\) on the
appropriate canonical encoding.  The family \(\mathcal C(G,s)\), the rank
\(r(G,s)\), intervention feasibility, \(\operatorname{LP}(G,s,c)\),
\(\operatorname{OPT}(G,s,c)\), and their approximation ratios are invariant
under causal-vertex relabeling; vertex-cover feasibility, optimum, and
approximation ratio are invariant under source-vertex relabeling.  Hence the
corresponding supplied order transports algorithms and certificates between
a generic finite carrier and its own canonical carrier with polynomial
overhead.  Elementary finite-graph, finite-set, comparison, and rational
arithmetic operations used in this paper have polynomially bounded
intermediate binary encodings and polynomial bit complexity.  Polynomial-time
maps on these encodings are closed under composition, including the
composition of private-guard expansion, parent-spouse rounding, verification,
and reverse guard replacement.

%%% FIELD proof-lem-fixed-point-hedge-oracle
For \(W\subseteq F\) with \(s\in W\), define
\[
 \mathsf A(W)=\{v\in W:v\text{ reaches }s\text{ in }D[W]\}
\]
and let \(\mathsf K(W)\) be the vertex set of the connected component of
\(s\) in \(B[W]\).  Starting from \(W_0=F\), alternate
\(W_{2j+1}=\mathsf A(W_{2j})\) and
\(W_{2j+2}=\mathsf K(W_{2j+1})\).  Both maps are deflationary and preserve
\(s\).  Each nonstationary step deletes a vertex.  Therefore, after at most
\(2|V|+2\) searches, a full two-step pass deletes nothing and the sequence
reaches a common fixed point \(F^{\star}\) satisfying
\[
 \mathsf A(F^{\star})=\mathsf K(F^{\star})=F^{\star}.
\]

If \(J\subseteq W\) is a hedge, every directed path to \(s\) inside
\(D[J]\) is also present in \(D[W]\), so \(J\subseteq\mathsf A(W)\).
Likewise every bidirected path inside the connected graph \(B[J]\) lies in
\(B[W]\), so \(J\subseteq\mathsf K(W)\).  Induction over the alternating
restrictions shows that every hedge contained in \(F\) is contained in
\(F^{\star}\).

At the common fixed point, every vertex reaches \(s\) in
\(D[F^{\star}]\), and \(B[F^{\star}]\) is connected.  If
\(F^{\star}\neq\{s\}\), Definition
\(\mathrm{def{:}singleton\mbox{-}hedges}\) therefore gives
\(F^{\star}\in\mathsf{Hdg}(G,s)\).  If
\(F^{\star}=\{s\}\), the survival implication proves that the original
\(F\) contained no hedge.  This establishes the equivalence.

Each restriction is one directed-reachability or undirected-connectivity
search on a canonically encoded induced causal graph.  The elementary
bit-cost clause of Assumption
\(\mathrm{ass{:}uniform\mbox{-}poly\mbox{-}time\mbox{-}model}\), together
with the linear bound on the number of searches, proves polynomial runtime
on the causal \(\mathrm{Fin}\,|V|\) carrier.  The causal relabeling clause,
using \(\triangleleft_V\), transfers the fixed set and hedge predicate to
every finite supplied-ordered causal carrier.

%%% FIELD proof-lem-rational-lp-polynomial-time
Lemma
\(\mathrm{lem{:}rational\mbox{-}lp\mbox{-}polynomial\mbox{-}time\mbox{-}source}\)
supplies an algorithm polynomial in the total binary length of a rational
linear program and guarantees an optimal basic solution whenever an optimum
exists.  For a canonically indexed program of dimensions \(d\) and total
binary length \(L\), its ordinary rational input length is bounded by a
polynomial in \(d\) and \(L\), and conversely these quantities are bounded
by its explicit encoding length.  The bit-cost clause of Assumption
\(\mathrm{ass{:}uniform\mbox{-}poly\mbox{-}time\mbox{-}model}\) charges
comparisons and rational field operations at their binary bit complexity.
Substituting the canonical encoding into the cited algorithm therefore gives
runtime polynomial in the dimensions and \(L\), with the same
optimal-basic-solution guarantee.  When the indexed program is obtained from
a causal instance, the causal relabeling clause transports its solution
along \(V\simeq\mathrm{Fin}\,|V|\) without changing feasibility or objective
value.

%%% FIELD proof-lem-ugc-vertex-cover-hardness
Assume \(\mathsf{UGC}\) and fix
\(\epsilon\in\mathbb R_{>0}\).  Lemma
\(\mathrm{lem{:}ugc\mbox{-}vertex\mbox{-}cover\mbox{-}hardness\mbox{-}source}\)
gives the published \((2-\epsilon)\)-inapproximability theorem for Minimum
Vertex Cover.  Suppose, contrary to the claimed canonical specialization,
that an algorithm with that ratio ran in polynomial time on canonically
\(\mathrm{Fin}\,|U|\)-encoded source graphs.  Given any finite source graph
equipped with \(\prec\), the source-carrier clause of Assumption
\(\mathrm{ass{:}uniform\mbox{-}poly\mbox{-}time\mbox{-}model}\) computes
its canonical adjacency-bit encoding and transports the returned vertex set
back along \(U\simeq\mathrm{Fin}\,|U|\) with polynomial overhead.
Vertex-cover feasibility, cardinality, optimum, and approximation ratio are
invariant under this source relabeling.  The transported output would
therefore be a polynomial-time \((2-\epsilon)\)-approximation on the source
theorem's finite graph instances, a contradiction.  Thus the published
lower bound holds on the declared canonical source carrier.  The same source
relabeling clause gives the promised transfer to a generic finite \(U\) in
later consumers; no causal carrier or causal order is used in this step.

%%% FIELD proof-thm-objective-preserving-hardness
Assign \(c(v)=1\) to every \(v\in V^{-}\).  Theorem
\(\mathrm{thm{:}private\mbox{-}guard\mbox{-}realization}\) gives
\[
 \mathcal C(G_H,s)
 =\bigl\{\{x_u,x_v,z_e\}:e=\{u,v\}\in E\bigr\}.
\]
Therefore the feasibility constraint indexed by \(e=\{u,v\}\) is exactly
\[
 A\cap\{x_u,x_v,z_e\}\neq\varnothing.
\]
By Definition \(\mathrm{def{:}mcid\mbox{-}objective}\), whose feasible
class is supplied by the cited hedge-hitting identification bridge, these
are precisely the proxy interventions that identify the causal target
\(Q_s\).

If \(C\subseteq U\) is a vertex cover, put
\(A_C=\{x_u:u\in C\}\).  For every edge \(e=\{u,v\}\), at least one of
\(u,v\) belongs to \(C\), so \(A_C\) meets the corresponding residual.
Thus \(A_C\) is feasible and
\[
 c(A_C)=|C|,
 \qquad
 \operatorname{OPT}(G_H,s,c)\leq\tau(H).
\]

Conversely, let \(A\subseteq V^{-}\) be feasible and initialize
\(C_0=\{u\in U:x_u\in A\}\).  Process the edges in the order induced by
\(\prec\).  If \(e=\{u,v\}\) is not covered by the current set, insert one
endpoint, say the \(\prec\)-smaller one.  At such a step neither \(x_u\)
nor \(x_v\) is in \(A\), so feasibility of the displayed residual
constraint forces \(z_e\in A\).  Charge the insertion to this private
guard.  Different edges have different guards, and an endpoint already
inserted creates no new charge.  The final set \(C_A\) is a vertex cover and
satisfies
\[
 |C_A|
 \leq |\{u:x_u\in A\}|+|\{e:z_e\in A\}|
 =|A|.
\]
The equality uses the fact that the manipulable vertices of \(G_H\) are
exactly its source copies and private guards.  Applying the map to an
optimal intervention gives
\[
 \tau(H)\leq |C_A|\leq\operatorname{OPT}(G_H,s,c).
\]
Together with the forward inequality this proves
\[
 \operatorname{OPT}(G_H,s,c)=\tau(H).
\]

We now certify the computational and approximation-preserving parts without
conflating carriers.  Canonically encode the source graph using the
isomorphism \(U\simeq\mathrm{Fin}\,|U|\) induced by \(\prec\).  The
expansion writes \(1+|U|+|E|\) vertices and the displayed adjacency entries.
Give its causal output the deterministic order that lists \(s\), then the
source copies in \(\prec\)-order, and then the guards in induced
lexicographic edge order; this yields the separate canonical causal encoding
on \(V(G_H)\).  The reverse map performs one ordered edge scan.  The two
canonicalization clauses and the elementary bit-cost clause of Assumption
\(\mathrm{ass{:}uniform\mbox{-}poly\mbox{-}time\mbox{-}model}\) make both
maps polynomial in their appropriate input lengths, and its composition
clause makes the guarded composition polynomial.  Its source relabeling
clause transfers source instances and covers along \(U\), while its causal
relabeling clause transfers interventions and the invariant clutter,
feasibility, and objective values along \(V\).

Equality of optima gives the first \(L\)-reduction inequality with
\(\alpha=1\).  Since \(C_A\) is a cover and \(|C_A|\leq|A|\),
\[
 0\leq |C_A|-\tau(H)
 \leq |A|-\operatorname{OPT}(G_H,s,c),
\]
which is the second inequality with \(\beta=1\).  Thus an exact MCID
algorithm would solve Vertex Cover, and any approximation error is carried
back with no enlargement.  Lemma
\(\mathrm{lem{:}cubic\mbox{-}vertex\mbox{-}cover\mbox{-}apx\mbox{-}complete}\)
states that Minimum Vertex Cover is \(\mathsf{APX}\)-complete, and hence
\(\mathsf{NP}\)-hard and \(\mathsf{APX}\)-hard, already at maximum degree
three.  The unit-constant \(L\)-reduction transfers both hardness
conclusions.  The private-guard theorem supplies rank three whenever
\(E\neq\varnothing\) and a bidirected tree of radius at most two.  Edgeless
source instances are recognized and solved by the empty set, so restricting
the hardness reduction to nonempty instances gives exact residual rank
three.  This proves every asserted restriction.

%%% FIELD proof-thm-rank-three-rounding
Fix the supplied causal-vertex order \(\triangleleft_V\) on \(V\).
Assumption \(\mathrm{ass{:}rank\mbox{-}three\mbox{-}promise}\) and
Definition \(\mathrm{def{:}residual\mbox{-}rank}\) give \(|R|\leq3\) for
every \(R\in\mathcal C(G,s)\).  Definitions
\(\mathrm{def{:}singleton\mbox{-}hedges}\) and
\(\mathrm{def{:}minimal\mbox{-}residual\mbox{-}clutter}\), together with
Theorem \(\mathrm{thm{:}first\mbox{-}rank\mbox{-}boundary}\), give
\(|R|\geq3\).  Thus each nonempty clutter member has size exactly three,
and the final clause of that theorem gives
\[
 R\subseteq P_s\cup T_s,
 \qquad R\cap P_s\neq\varnothing,
 \qquad R\cap T_s\neq\varnothing,
 \qquad P_s\cap T_s=\varnothing.
\]

The all-ones vector is feasible for the covering LP, and its feasible set
is a closed subset of the compact cube \([0,1]^{V^{-}}\); its linear
objective therefore attains a minimum.  Proposition
\(\mathrm{prop{:}promised\mbox{-}enumeration}\), at promised bound three
and order \(\triangleleft_V\), lists all constraints.  On the canonical
causal encoding of Assumption
\(\mathrm{ass{:}uniform\mbox{-}poly\mbox{-}time\mbox{-}model}\), the list
has polynomial length and rational entries of polynomial bit length.
Lemma \(\mathrm{lem{:}rational\mbox{-}lp\mbox{-}polynomial\mbox{-}time}\)
therefore returns an optimal rational vector \(y\).  Fix one such vector.

For \(0<t<1/2\), form \(A_t\) according to Definition
\(\mathrm{def{:}parent\mbox{-}spouse\mbox{-}rounding}\).  Suppose that it
misses some \(R\in\mathcal C(G,s)\), and let
\(a=|R\cap P_s|\).  The preceding partition gives \(a\in\{1,2\}\).
Missed parent and spouse vertices satisfy respectively
\[
 y_p<t,
 \qquad y_b<\frac12-t.
\]
If \(a=1\), then
\[
 \sum_{v\in R}y_v
 <t+2\left(\frac12-t\right)=1-t<1;
\]
if \(a=2\), then
\[
 \sum_{v\in R}y_v
 <2t+\left(\frac12-t\right)=\frac12+t<1.
\]
Both contradict the LP constraint \(\sum_{v\in R}y_v\geq1\).  Hence every
\(A_t\) hits the entire minimal residual clutter and is a feasible
identifying intervention by Definition
\(\mathrm{def{:}mcid\mbox{-}objective}\).

Let \(T\) be uniform on \((0,1/2)\); this is only an averaging device over
deterministic thresholds.  Interval lengths give
\[
 \Pr(p\in A_T)=\min\{1,2y_p\}\quad(p\in P_s),
 \qquad
 \Pr(b\in A_T)=\min\{1,2y_b\}\quad(b\in T_s).
\]
No other vertex is selected.  Disjointness of \(P_s,T_s\), linearity of
expectation, nonnegativity of \(y\), and positivity of \(c\) yield
\[
 \begin{aligned}
 \mathbb E\!\left[\sum_{v\in A_T}c(v)\right]
 &=\sum_{v\in P_s\cup T_s}c(v)\min\{1,2y_v\}\\
 &\leq2\sum_{v\in V^{-}}c(v)y_v
 =2\operatorname{LP}(G,s,c).
 \end{aligned}
\]
Every feasible integral intervention has a feasible incidence vector for
the LP with the same cost, so
\[
 \operatorname{LP}(G,s,c)\leq\operatorname{OPT}(G,s,c).
\]

Membership in \(A_t\) changes only at the finitely many critical values
\[
 \{y_p:p\in P_s\}\cup
 \left\{\frac12-y_b:b\in T_s\right\}
\]
that lie in \((0,1/2)\).  On each open cell between consecutive values the
set is constant; the critical values themselves have Lebesgue measure zero.
If every positive-length cell had cost greater than
\(2\operatorname{LP}(G,s,c)\), the expectation would too.  Hence some
cell, and therefore the cheapest candidate evaluated by Definition
\(\mathrm{def{:}parent\mbox{-}spouse\mbox{-}rounding}\), has cost at most
\(2\operatorname{LP}(G,s,c)\).  The rationality of \(y\) makes every
critical value and, for example, every cell midpoint rational.

There are at most \(2|V|\) critical values.  Promised enumeration, the
canonical rational-LP algorithm, rational sorting, midpoint construction,
and candidate evaluation each use polynomially many elementary operations.
The bit-cost and composition clauses of Assumption
\(\mathrm{ass{:}uniform\mbox{-}poly\mbox{-}time\mbox{-}model}\) make the
composite algorithm polynomial in \(|V|\) and the causal encoding length
\(L\).  Its causal relabeling clause, using \(\triangleleft_V\), transfers
the invariant result from \(\mathrm{Fin}\,|V|\) to the generic finite
causal carrier.  Thus the returned threshold satisfies
\[
 \sum_{v\in A_t}c(v)
 \leq2\operatorname{LP}(G,s,c)
 \leq2\operatorname{OPT}(G,s,c),
\]
as required.

%%% FIELD proof-prop-promised-enumeration
Fix a nonnegative integer \(r_0\), a supplied causal order
\(\triangleleft_V\) on \(V\), and the promised condition
\(r(G,s)\leq r_0\).  Generate once, in induced lexicographic order, every
\(R\subseteq V^{-}\) with \(|R|\leq r_0\), and put
\(F_R=R\cup\{s\}\).  The case \(R=\varnothing\) fails the strict
containment clause in Definition \(\mathrm{def{:}singleton\mbox{-}hedges}\).
For nonempty \(R\), a bidirected search in \(B[F_R]\) tests connectivity,
and a reverse directed search from \(s\) in \(D[F_R]\) tests whether every
vertex reaches \(s\).  Hence
\[
 F_R\in\mathsf{Hdg}(G,s)
 \quad\Longleftrightarrow\quad
 B[F_R]\text{ is connected and every }v\in F_R
 \text{ reaches }s\text{ in }D[F_R].
\]

For a candidate hedge, perform the same test on every
\(R'\subsetneq R\).  Every proper subset of \(F_R\) that contains \(s\)
is uniquely \(R'\cup\{s\}\), so Definition
\(\mathrm{def{:}minimal\mbox{-}residual\mbox{-}clutter}\) gives
\[
 R\in\mathcal C(G,s)
 \quad\Longleftrightarrow\quad
 F_R\in\mathsf{Hdg}(G,s)
 \text{ and }F_{R'}\notin\mathsf{Hdg}(G,s)
 \text{ for all }R'\subsetneq R.
\]
Every retained set is therefore a clutter member.  Conversely, the rank
promise and Definition \(\mathrm{def{:}residual\mbox{-}rank}\) put every
clutter member among the tested candidates.  This proves exact-once output
and the requested order.

Writing \(n=|V|\), there are
\[
 \sum_{j=0}^{r_0}\binom{n-1}{j}=O(n^{r_0})
\]
candidates for fixed \(r_0\), and at most \(2^{r_0}\) proper subsets per
candidate.  Each graph search is polynomial.  Candidates and outputs may be
streamed while storing only the current candidate, one proper subset, and
the graph-search workspace, proving polynomial space.  On the canonical
causal carrier, the causal-encoding and elementary bit-cost clauses of
Assumption \(\mathrm{ass{:}uniform\mbox{-}poly\mbox{-}time\mbox{-}model}\)
turn this operation count into time \(|V|^{r_0+O(1)}\), polynomial in the
remaining causal input length \(L\).  The causal relabeling clause, using
\(\triangleleft_V\), transfers the invariant output to the generic finite
causal carrier.

At \(r_0=3\), Assumption
\(\mathrm{ass{:}rank\mbox{-}three\mbox{-}promise}\) supplies the promise.
Theorem \(\mathrm{thm{:}first\mbox{-}rank\mbox{-}boundary}\) gives the
matching lower bound three on nonempty residuals.  Definition
\(\mathrm{def{:}promised\mbox{-}enumerator}\) therefore tests precisely
the required triples, of which there are \(O(n^3)\).

For the scoped exact search, use this enumerated rank-three clutter.  At a
state \((A_0,j)\), return \(A_0\) if it hits every residual.  If it does
not and \(j=0\), reject the branch.  Otherwise select an uncovered
three-set \(R\) and recurse on
\[
 (A_0\cup\{v\},j-1),
 \qquad v\in R,
\]
retaining a least-cost successful result.  Induction on \(j\) proves
optimality among feasible extensions using at most \(j\) new vertices.
The base case is the stated feasibility test.  In the inductive step, if
\(A_0\) is feasible then positivity of \(c\) makes every strict extension
more expensive.  If it is infeasible, every feasible extension must hit
the selected uncovered \(R\), hence occurs in one of the three branches;
the induction hypothesis optimizes each branch.

The recursion tree has at most
\[
 1+3+\cdots+3^k=O(3^k)
\]
nodes.  Every node performs polynomial work on the polynomial-size
enumerated clutter.  Rational sums and comparisons have bit length
polynomial in \(|V|\) and \(L\).  The elementary bit-cost and composition
clauses of Assumption
\(\mathrm{ass{:}uniform\mbox{-}poly\mbox{-}time\mbox{-}model}\) yield the
claimed \(O(3^k\operatorname{poly}(|V|,L))\) runtime on the canonical
causal carrier, and its causal relabeling clause transfers the result along
\(V\simeq\mathrm{Fin}\,|V|\).  The proof uses the explicit cardinality
budget \(k\) and makes no unconstrained weighted \(3^k\) claim.

%%% FIELD proof-prop-optimizer-or-violation
Fix the supplied causal order \(\triangleleft_V\) on \(V\) and define the
proof-local family
\[
 \mathcal C_{\leq3}=\{R\in\mathcal C(G,s):|R|\leq3\}.
\]
Run Definition \(\mathrm{def{:}promised\mbox{-}enumerator}\) on all
subsets of size at most three.  The connectivity, reachability, and
proper-subset tests are exactly Definitions
\(\mathrm{def{:}singleton\mbox{-}hedges}\) and
\(\mathrm{def{:}minimal\mbox{-}residual\mbox{-}clutter}\), so the output
is exactly \(\mathcal C_{\leq3}\), without a global rank promise.  It has
polynomial size.  Theorem
\(\mathrm{thm{:}first\mbox{-}rank\mbox{-}boundary}\) says that every
member is a three-set contained in the disjoint union
\(P_s\mathbin{\dot\cup}T_s\) and meeting both parts.

Define the restricted values
\[
 \operatorname{LP}_{\leq3}
 =\min\left\{\sum_{v\in V^{-}}c(v)y_v:
 y\in[0,1]^{V^{-}},\ \sum_{v\in R}y_v\geq1
 \text{ for every }R\in\mathcal C_{\leq3}\right\}
\]
and
\[
 \operatorname{OPT}_{\leq3}
 =\min\left\{\sum_{v\in A}c(v):
 A\subseteq V^{-},\ A\cap R\neq\varnothing
 \text{ for every }R\in\mathcal C_{\leq3}\right\}.
\]
These are proof-local restrictions of the frozen global functionals.  The
restricted LP has rational input and a nonempty compact feasible region.
On the canonical causal encoding, Lemma
\(\mathrm{lem{:}rational\mbox{-}lp\mbox{-}polynomial\mbox{-}time}\)
returns an optimal rational vector \(y^{(3)}\).

Apply the parent-spouse threshold formula to \(y^{(3)}\).  If a threshold
set missed a residual with one parent, then
\[
 \sum_{v\in R}y^{(3)}_v
 <t+2\left(\frac12-t\right)=1-t<1;
\]
if it missed one with two parents, then
\[
 \sum_{v\in R}y^{(3)}_v
 <2t+\left(\frac12-t\right)=\frac12+t<1.
\]
Both contradict the restricted LP.  Thus every threshold set hits
\(\mathcal C_{\leq3}\).  With \(T\) uniform on \((0,1/2)\), interval
length gives
\[
 \Pr(p\in A_T)=\min\{1,2y^{(3)}_p\},
 \qquad
 \Pr(b\in A_T)=\min\{1,2y^{(3)}_b\}.
\]
Linearity, positivity of costs, and disjointness of the two parts imply
\[
 \mathbb E[c(A_T)]\leq2\operatorname{LP}_{\leq3}.
\]
The rational-breakpoint procedure therefore returns
\(A^{\mathrm{rnd}}\) with
\[
 c(A^{\mathrm{rnd}})\leq2\operatorname{LP}_{\leq3}
 \leq2\operatorname{OPT}(G,s,c),
\]
where the last inequality holds because every globally feasible incidence
vector is feasible for the restricted LP.  Alternatively, an exact
optimization backend may return an integral minimizer \(A^{(3)}\) with
\(c(A^{(3)})=\operatorname{OPT}_{\leq3}\); no polynomial-time claim is
made for this exact integral alternative.

For either returned \(A\), put \(F_0=V\setminus A\), which contains \(s\),
and apply Lemma
\(\mathrm{lem{:}fixed\mbox{-}point\mbox{-}hedge\mbox{-}oracle}\).  If the
fixed point is \(\{s\}\), then \(F_0\) contains no hedge.  Since the graph
is finite, every hedge contains an inclusion-minimal hedge; hence \(A\)
hits all of \(\mathcal C(G,s)\).  Definition
\(\mathrm{def{:}mcid\mbox{-}objective}\) then certifies global feasibility
and identification of \(Q_s\).  For \(A^{(3)}\), restricted feasibility
gives \(\operatorname{OPT}_{\leq3}\leq\operatorname{OPT}(G,s,c)\), while
successful global feasibility gives
\[
 \operatorname{OPT}(G,s,c)
 \leq c(A^{(3)})=\operatorname{OPT}_{\leq3};
\]
thus the exact restricted optimum is globally exact.  For
\(A^{\mathrm{rnd}}\), the already established inequality gives the global
factor-two certificate.

If verification fails, the oracle returns a hedge \(F\subseteq F_0\).
For each \(v\in F\setminus\{s\}\), run the oracle on
\(F\setminus\{v\}\).  Whenever a hedge survives, replace \(F\) by the
strictly smaller fixed-point hedge returned and restart.  At most
\(|V|-1\) replacements occur, with at most \(|V|-1\) oracle calls between
replacements.  At termination no single deletion contains a hedge.  If a
proper hedge \(J\subsetneq F\) existed, choose
\(v\in F\setminus J\); both hedges contain \(s\), so \(v\neq s\), and
\(J\subseteq F\setminus\{v\}\), a contradiction.  Therefore
\[
 R=F\setminus\{s\}\in\mathcal C(G,s).
\]
Since \(F\subseteq V\setminus A\), it is missed by \(A\); because both
candidate types hit every member of \(\mathcal C_{\leq3}\), necessarily
\(|R|>3\).  The final graph searches verify the hedge predicates and the
failed single-deletion searches verify minimality.

Canonical bounded enumeration, rational LP solution and rounding,
fixed-point verification, and deletion minimization perform polynomially
many elementary operations on the causal encoding.  The elementary bit-cost
and composition clauses of Assumption
\(\mathrm{ass{:}uniform\mbox{-}poly\mbox{-}time\mbox{-}model}\) make their
composition polynomial on \(\mathrm{Fin}\,|V|\).  Its causal relabeling
clause, using \(\triangleleft_V\), transfers the invariant certificate and
guarantees to the generic finite causal carrier.  No source carrier or
source order is used by this interface.  This proves the assertion without
recognizing the rank promise.

%%% FIELD proof-thm-tight-frontier
Consider any instance in the stated class and choose
\(R_0\in\mathcal C(G,s)\).  It is nonempty by Definitions
\(\mathrm{def{:}singleton\mbox{-}hedges}\) and
\(\mathrm{def{:}minimal\mbox{-}residual\mbox{-}clutter}\).  Positivity of
the rational costs and finiteness give
\(m_0=\min_{v\in R_0}c(v)>0\).  Every feasible LP vector satisfies
\[
 \sum_{v\in V^{-}}c(v)y_v
 \geq\sum_{v\in R_0}c(v)y_v
 \geq m_0\sum_{v\in R_0}y_v
 \geq m_0>0.
\]
Hence \(\operatorname{LP}(G,s,c)>0\).  Choose a causal-vertex order
\(\triangleleft_V\) on the finite set \(V\) and invoke Theorem
\(\mathrm{thm{:}rank\mbox{-}three\mbox{-}rounding}\).  It gives a feasible
identifying intervention \(A_t\) with
\[
 \operatorname{OPT}(G,s,c)
 \leq\sum_{v\in A_t}c(v)
 \leq2\operatorname{LP}(G,s,c).
\]
Division by the positive LP value proves the universal ratio bound two.

For the matching sequence, let \(n\geq2\), equip \(K_n\) with any source
order \(\prec\) on \(U\), form \(G_{K_n}\), and give every manipulable
vertex unit cost.  Theorems
\(\mathrm{thm{:}private\mbox{-}guard\mbox{-}realization}\) and
\(\mathrm{thm{:}objective\mbox{-}preserving\mbox{-}hardness}\) give a
radius-two bidirected tree, residual rank three, and
\[
 \mathcal C(G_{K_n},s)
 =\bigl\{\{x_u,x_v,z_e\}:e=\{u,v\}\in E(K_n)\bigr\},
 \qquad
 \operatorname{OPT}(G_{K_n},s,c)=\tau(K_n)=n-1.
\]
The last equality holds because omitting two vertices misses their edge,
whereas all but one vertex cover every edge.

Set \(y_{x_u}=1/2\) for every \(u\) and \(y_{z_e}=0\) for every \(e\).
Each clutter constraint has sum one, so this solution is feasible and has
cost \(n/2\).  Conversely, summing all clutter constraints gives
\[
 (n-1)\sum_{u\in U}y_{x_u}+\sum_{e\in E(K_n)}y_{z_e}
 \geq\binom n2.
\]
After division by \(n-1>0\),
\[
 \sum_{u\in U}y_{x_u}
 +\frac1{n-1}\sum_{e\in E(K_n)}y_{z_e}\geq\frac n2.
\]
Since every coordinate is nonnegative and \(1/(n-1)\leq1\), the unit-cost
LP objective is at least the left side.  Therefore
\[
 \operatorname{LP}(G_{K_n},s,c)=\frac n2,
 \qquad
 \frac{\operatorname{OPT}(G_{K_n},s,c)}
      {\operatorname{LP}(G_{K_n},s,c)}
 =2-\frac2n.
\]
These ratios tend to two, proving that the supremum equals two.

Finally assume \(\mathsf{UGC}\), fix the positive real
\(\epsilon\), and suppose that a polynomial-time
\((2-\epsilon)\)-approximation exists on the narrower causal subclass.
Start with a source graph \(H=(U,E)\) equipped with its independent source
order \(\prec\).  Canonicalize it through
\(U\simeq\mathrm{Fin}\,|U|\), form the private-guard graph, equip that
causal output with the deterministic causal order used in the hardness
proof, and apply the hypothetical algorithm.  If \(E\neq\varnothing\),
the private-guard theorem puts the output in the claimed subclass.  The
reverse map of Theorem
\(\mathrm{thm{:}objective\mbox{-}preserving\mbox{-}hardness}\) returns a
cover \(C_A\) satisfying
\[
 |C_A|\leq|A|
 \leq(2-\epsilon)\operatorname{OPT}(G_H,s,c)
 =(2-\epsilon)\tau(H).
\]
If \(E=\varnothing\), return the empty cover.  The two canonicalization,
relabeling, elementary bit-cost, and guarded-composition clauses of
Assumption \(\mathrm{ass{:}uniform\mbox{-}poly\mbox{-}time\mbox{-}model}\)
make this a polynomial-time approximation on the canonical source carrier
and transfer it to generic finite \(U\); they never identify \(U\) with the
distinct causal carrier \(V(G_H)\).  This contradicts the canonically scoped
Lemma \(\mathrm{lem{:}ugc\mbox{-}vertex\mbox{-}cover\mbox{-}hardness}\).
Thus no such approximation exists.
